%%%%%%%%%%%%%%%%%%%%%%%%%%%%%%%%%%%%%%%%%%%%%%%%%%%%%%%%%%%%%%%%%%%%%%
% Lớp 8 | Chương I. Đa thức | Đơn thức
% Dạng 8D1?1-2: Đơn thức đồng dạng; cộng, trừ đơn thức đồng dạng
%%%%%%%%%%%%%%%%%%%%%%%%%%%%%%%%%%%%%%%%%%%%%%%%%%%%%%%%%%%%%%%%%%%%%%


% nguồn: De_AI_Toan8_cuoi_hoc_ki_i_1 (nhập 21/11/2026)
\begin{bt}%[8D1N1-2]%[ma:b865d925]
	Sắp xếp các đơn thức sau thành từng nhóm các đơn thức đồng dạng:
$5x^2y$; $-3xy^2$; $-\dfrac{1}{2}x^2y$; $7xy^2$; $x^2y$; $-xy^2$; $4xyz$.
	\loigiai{
		Các nhóm đơn thức đồng dạng (có hệ số khác $0$ và cùng phần biến) là: \\
- Nhóm 1 (phần biến $x^2y$): $5x^2y$; $-\dfrac{1}{2}x^2y$; $x^2y$. \\
- Nhóm 2 (phần biến $xy^2$): $-3xy^2$; $7xy^2$; $-xy^2$. \\
- Đơn thức $4xyz$ không đồng dạng với đơn thức nào trong danh sách trên.
	}
\end{bt}

% nguồn: De_AI_Toan8_cuoi_hoc_ki_i_1 (nhập 21/11/2026)
\begin{bt}%[8D1N1-2]%[ma:6f407024]
	Tính tổng, hiệu các đơn thức đồng dạng sau:
a) $A = 4x^2y + 7x^2y - 3x^2y$;
b) $B = 2xy^3 - 5xy^3 + \dfrac{1}{2}xy^3$.
	\loigiai{
		a) $A = (4 + 7 - 3)x^2y = 8x^2y$. \\
b) $B = \left(2 - 5 + \dfrac{1}{2}\right)xy^3 = -\dfrac{5}{2}xy^3$.
	}
\end{bt}

% nguồn: De_AI_Toan8_cuoi_hoc_ki_i_1 (nhập 21/11/2026)
\begin{bt}%[8D1N1-2]%[ma:0f20e080]
	Thực hiện phép tính:
a) $C = -3x^3y^2 + 8x^3y^2 - 5x^3y^2$;
b) $D = \dfrac{1}{3}xyz + \dfrac{2}{3}xyz - 4xyz$.
	\loigiai{
		a) $C = (-3 + 8 - 5)x^3y^2 = 0 \cdot x^3y^2 = 0$. \\
b) $D = \left(\dfrac{1}{3} + \dfrac{2}{3} - 4\right)xyz = (1 - 4)xyz = -3xyz$.
	}
\end{bt}

% nguồn: De_AI_Toan8_cuoi_hoc_ki_i_1 (nhập 21/11/2026)
\begin{bt}%[8D1N1-2]%[ma:dd15e15b]
	Thu gọn biểu thức $E = 6x^2yz - 2xyzx + 5zx^2y$ rồi xác định hệ số và bậc của đơn thức thu được.
	\loigiai{
		Ta thu gọn từng hạng tử: $2xyzx = 2x^2yz$ và $5zx^2y = 5x^2yz$. \\
Do đó $E = 6x^2yz - 2x^2yz + 5x^2yz = (6 - 2 + 5)x^2yz = 9x^2yz$. \\
Đơn thức $E$ có hệ số là $9$ và bậc là $2 + 1 + 1 = 4$.
	}
\end{bt}

% nguồn: De_AI_Toan8_cuoi_hoc_ki_i_1 (nhập 21/11/2026)
\begin{bt}%[8D1N1-2]%[ma:ae1dbff9]
	Tìm số tự nhiên $n$ để hai đơn thức $2x^{n+1}y^4$ và $5x^3y^{n+2}$ đồng dạng. Khi đó hãy tính tổng của hai đơn thức.
	\loigiai{
		Để hai đơn thức đồng dạng thì số mũ của từng biến tương ứng phải bằng nhau: \\
$n + 1 = 3$ và $n + 2 = 4$, suy ra $n = 2$ (thỏa mãn $n \in \mathbb{N}$). \\
Khi $n = 2$, hai đơn thức là $2x^3y^4$ và $5x^3y^4$. Tổng của chúng là $2x^3y^4 + 5x^3y^4 = 7x^3y^4$.
	}
\end{bt}
